\begin{center}
[4].\\
\textsc{Gauss} to \textsc{Olbers}.\quad Göttingen, 13~January 1821.
\end{center}

\par\noindent\hspace*{\parindent}.\ .\ .\ .\ .\ .\quad If all the external circumstances turn out to be more favourable during the actual survey than they have appeared so far, I believe I could well enjoy the work, and then I would also be very glad to undertake an extension of the triangulation to the west if I were charged with it.
A connection with \textsc{Krayenhoff}'s triangles is certainly desirable, but where are they to be found?
I do not know whether they have been printed anywhere, and my poor success with \textsc{Epailly}'s triangles leaves me quite discouraged.
Even \textsc{Laplace}, to whom I
\par\clearpage
\noindent368\hfill{\small\textbf{CORRESPONDENCE WITH OLBERS.}}\par\medskip
\noindent wrote about 9 weeks ago, has not replied at all.
In my opinion, all well-measured triangles of the 1st order should be regarded as something to which the entire public has a claim, and the whole of Europe should gradually be covered with such triangles.
Some years ago I worked out a method of my own for handling such measurements as efficiently as possible, since I find everything I have read about it utterly worthless.
For example, many mathematicians have gone to great lengths over the problem of computing longitude and latitude from distances from the meridian and the perpendicular, taking into account the elliptical shape of the Earth, while, as far as I know, no one has first asked

\par\noindent\hspace*{\parindent}1.\ how those distances, understood as they usually are, can be obtained from the survey with equal precision, since it appears that most people perform this computation as if in the plane or give altogether incorrect or unusable rules for it,

\par\noindent\hspace*{\parindent}2.\ whether it is expedient at all to use \textls[120]{distances understood in this way}, since it is certain that deriving them from the triangles with sufficient precision can be done only by extremely laborious calculations, just as passing back from them to the longitudes and latitudes also takes considerable effort.
The whole thing would simply be “putting the cart before the horse.”
“If anything useful is to be placed \textls[120]{between} the triangles and the longitudes and latitudes, it must be something quite different from those coordinates as ordinarily understood.”
I cannot explain in detail here how this is done in my theory. I will only remark that what I put between the triangles and the longitudes and latitudes \textls[120]{are the coordinates 1)~with which every point can most conveniently be represented in a plane}.
These coordinates are obtained very readily and easily from the measured triangles, without requiring very precise knowledge of the flattening of the Earth, and 2)~the longitude and latitude in turn follow from them just as easily, although one must of course know the flattening.
I intend to make this theory known with my future survey results, if I do not do so sooner, and for now I ask you to keep these ideas I have sketched to yourself.
I would very much like to apply them to the Hanover triangles and to all the others that can be connected with them
\par\clearpage
\noindent{\small\textbf{ON THE HANOVER TRIANGULATION.}}\hfill369\par\medskip
\noindent and thus give a \textls[120]{Description géométrique} of a large part of Europe, if I were given adequate support through the sharing of results.
But!!

\par\noindent\hspace*{\parindent}Mr.~v.\ \textsc{Müffling} did send me his 15 triangles from the Rhine to Seeberg.
I have already connected Göttingen provisionally, though of course only very roughly.
For in 1812 I measured, on Hanstein, the angles between Göttingen, Brocken, and the Boineburg (the 2nd \textsc{Müffling} point). The position of Hanstein relative to Göttingen follows approximately from my angle measurements around here. The angles are admittedly uncertain by several minutes, but I believe the resulting placement of Göttingen is probably reliable to within about 100 metres.
The resulting difference in longitude between the Göttingen and Seeberg observatories is \mbox{3\textsuperscript{m}8\rlap{\textsuperscript{s}},7} in time, which agrees very closely with the astronomical determinations.
Thus, if Paris is \mbox{33\textsuperscript{m}35\textsuperscript{s}} west of Seeberg, it would be \mbox{30\textsuperscript{m}26\rlap{\textsuperscript{s}},3} west of Göttingen.
The new observatory lies \mbox{1\rlap{\textsuperscript{s}},9} east of the old one, which I formerly always placed \mbox{30\textsuperscript{m}23\(\tfrac{1}{2}\)\textsuperscript{s}} from Paris.
